\section{10차시 --- LED 5개 시퀀스와 카운트다운}
\label{sec:ch10}

본 차시는 LED를 5개로 확장하여 \emph{발사 시퀀스}를 구현한다. 순차
점등, 파티 조명, 카운트다운(역순 소등), 그리고 이들을 통합한 조명쇼로
이어진다. 이 차시부터 학습자는 "하나의 시나리오를 코드 블록의 조합으로
구성한다"는 통합적 사고를 시작한다.

\subsection{미션 1 --- 발사 준비 순차 점등!}
\label{subsec:ch10-m1}
\missionmeta{LED 5개}{SEQUENCE}

\begin{storybox}
\sayares{알비! 발사 전에 LED가 하나씩 순서대로 켜지는 발사 준비 표시를 만들어보자!}
\sayalbi{삐릿- 좋아요! 1번부터 5번까지 순서대로 켜지면 발사 준비 단계가 완료되는 거예요!}
\end{storybox}

\subsubsection*{학습 목표}
\begin{itemize}[leftmargin=1.4em]
  \item LED1→LED2→LED3→LED4→LED5 순서대로 0.5초 간격 켜기
  \item 5개 모두 켜진 후 3번 동시 깜빡임
  \item 순차 점등으로 발사 준비 카운트다운 완성
\end{itemize}

\begin{labbox}{샘플 코드 --- 발사 준비 순차 점등}
\begin{minted}{python}
# 발사 준비 순차 점등
for i in range(1, 6):
    led_on(i, 1.0)
    time.sleep(0.5)
# 준비 완료 신호
for j in range(3):
    for i in range(1,6): led_off(i)
    time.sleep(0.3)
    for i in range(1,6): led_on(i,1.0)
    time.sleep(0.3)
\end{minted}
\end{labbox}

\subsection{미션 2 --- 발사 파티 조명!}
\label{subsec:ch10-m2}
\missionmeta{LED 5개}{SEQUENCE}

\begin{storybox}
\sayares{발사 성공 후에는 파티 조명을 켜야지! LED 5개가 신나게 깜빡이는 파티 분위기!}
\sayalbi{삐릿♪ 야호! 발사 성공 파티예요! 화려하게 깜빡이는 파티 조명을 만들어봐요!}
\end{storybox}

\subsubsection*{학습 목표}
\begin{itemize}[leftmargin=1.4em]
  \item 5개 LED 동시에 빠르게 10번 깜빡이기
  \item 파도처럼 흐르는 패턴 (1→2→3→4→5→4→3→2→1)
  \item 세 가지 패턴 조합한 파티 조명 완성
\end{itemize}

\begin{labbox}{샘플 코드 --- 파티 조명}
\begin{minted}{python}
# 파티 조명!
for j in range(10):
    for i in range(1,6): led_on(i,1.0)
    time.sleep(0.2)
    for i in range(1,6): led_off(i)
    time.sleep(0.2)
# 파도 패턴
for rep in range(3):
    for i in range(1,6):
        led_on(i,1.0); time.sleep(0.15); led_off(i)
\end{minted}
\end{labbox}

\subsection{미션 3 --- LED 카운트다운 5→1!}
\label{subsec:ch10-m3}
\missionmeta{LED 5개}{SEQUENCE}

\begin{storybox}
\sayares{알비! LED 5개로 숫자 카운트다운을 표현할 수 있어? 5개가 켜진 상태에서 하나씩 꺼지면 5·4·3·2·1이 되잖아!}
\sayalbi{삐릿- 오, 천재적인 아이디어예요! LED가 하나씩 꺼지면서 발사 카운트다운을 표현하는 거군요!}
\end{storybox}

\subsubsection*{학습 목표}
\begin{itemize}[leftmargin=1.4em]
  \item LED 5개 모두 켜기
  \item 1초마다 하나씩 끄기 (5→4→3→2→1→모두 꺼짐)
  \item 모두 꺼진 순간 전체 3번 빠르게 깜빡임 (발사!)
\end{itemize}

\begin{labbox}{샘플 코드 --- 카운트다운 소등}
\begin{minted}{python}
# 카운트다운 소등
for i in range(1,6): led_on(i,1.0)
for i in range(5, 0, -1):
    led_off(i)
    time.sleep(1)
# 발사!
for j in range(3):
    for i in range(1,6): led_on(i,1.0)
    time.sleep(0.2)
    for i in range(1,6): led_off(i)
    time.sleep(0.2)
\end{minted}
\end{labbox}

\subsection{미션 4 --- 나만의 발사 조명쇼!}
\label{subsec:ch10-m4}
\missionmeta{LED 5개}{SEQUENCE}

\begin{storybox}
\sayares{알비, 지금까지 만든 것들을 합쳐서 완전한 발사 조명쇼를 만들어보자!}
\sayalbi{삐릿♪ 순차 점등 → 카운트다운 → 발사 → 파티 순서로 하나의 이야기처럼 연결!}
\end{storybox}

\subsubsection*{학습 목표}
\begin{itemize}[leftmargin=1.4em]
  \item 순차 점등(발사 준비) → 카운트다운(5→1)
  \item 발사 깜빡임 → 파티 조명
  \item 전체 흐름이 하나로 이어지는 조명쇼
\end{itemize}

\begin{labbox}{샘플 코드 --- 통합 발사 조명쇼}
\begin{minted}{python}
# 통합 발사 조명쇼
# 1. 발사 준비
for i in range(1,6): led_on(i,1.0); time.sleep(0.3)
# 2. 카운트다운
for i in range(5,0,-1): led_off(i); time.sleep(1)
# 3. 발사!
for j in range(5):
    for i in range(1,6): led_on(i,1.0)
    time.sleep(0.15)
    for i in range(1,6): led_off(i)
    time.sleep(0.15)
# 4. 파티!
for rep in range(3):
    for i in range(1,6): led_on(i,1.0); time.sleep(0.1); led_off(i)
\end{minted}
\end{labbox}

\begin{keypoint}
미션 4는 "서브시퀀스의 합성"을 명시적으로 가르치는 첫 미션이다.
앞선 3개 미션의 코드 패턴을 그대로 "복사 + 이어붙이기"하면 통합
조명쇼가 완성된다는 점을 강조하여, 학습자가 \emph{모듈화 사고}에
한 발 내디딜 수 있게 한다.
\end{keypoint}
